\chapter{Canonical-form reductions after blocking}\label{ch:canonical_after_blocking}

\section{Scope}

This chapter continues the canonical-form reduction begun in
Chapter~\ref{ch:canonical}. There the per-tensor steps were established:
invariant-subspace splitting, irreducible block decomposition, separation of the
all-zero summand, blockwise trace-preserving gauge, and one-block period removal.
Here the same chain is carried past a common blocking length, producing the
primitive-block decomposition that closes the reduction.

The chapter opens with the formal zero-block separation
(Theorem~\ref{thm:zero_block_separation}), which is the statement used by the
arbitrary-tensor TP gauge of Chapter~\ref{ch:canonical}
(Theorem~\ref{thm:tp_gauge_arbitrary}). It then assembles the after-blocking
decompositions and reaches two final statements: the primitive-block
decomposition of a single tensor (Theorem~\ref{thm:tp_primitive_blockdecomp})
and the joint decomposition of two same-MPV tensors at a common blocking period
(Theorem~\ref{thm:ft_after_blocking_common_length_common_sector_theorem}). Both
feed the proof of the Fundamental Theorem in Chapter~\ref{ch:ft_proof}.

\section{Zero-block separation}%
\label{sec:zero_block_separation}

In an irreducible block decomposition some blocks may have all-zero Kraus
operators. Such blocks contribute only at length $N=0$, where the contribution
equals their bond dimension. There is no zero tensor in the output: the
contribution is isolated as a single bond-dimension count $D_0$ in the
length-zero identity $D = D_0 + \sum_k D_k$, while the positive-length content is
an equality of MPV families.

\begin{lemma}[All-zero irreducible blocks have bond dimension at most one]%
    \label{lem:irreducible_allZero_dim_le_one}
    \lean{MPSTensor.isIrreducibleTensor_allZero_dim_le_one}
    \leanok
    \uses{def:irreducible_tensor}
    If $A$ is irreducible with every Kraus operator $A^i = 0$, then $D \le 1$.
\end{lemma}

\begin{proof}\leanok
    If $D \ge 2$, irreducibility forces some Kraus operator to be nonzero;
    otherwise the bond space splits as a direct sum of any two complementary
    subspaces. Hence $A^i = 0$ for every $i$ implies $D \le 1$.
\end{proof}

\begin{theorem}[Zero-block separation]%
    \label{thm:zero_block_separation}
    \lean{MPSTensor.exists_irreducible_blockDecomp_nonzeroBlocks}
    \leanok
    \uses{def:irreducible_tensor, def:same_mpv2_pos}
    Every MPS tensor $A$ admits a block decomposition into a zero-block bond
    dimension $D_0$ and a finite family of nonzero blocks $(B_k)_{k=1}^r$, each
    irreducible, with at least one nonzero Kraus operator, and of positive bond
    dimension, recorded by two facts. The positive-length equality of MPV
    families: for every $N \ge 1$ and every $\sigma$,
    \[
        V^{(N)}(A)_\sigma
        = V^{(N)}\!(\textstyle\bigoplus_{k=1}^r B_k)_\sigma.
    \]
    The length-zero dimension identity
    \[
        D = D_0 + \sum_{k=1}^r D_k,
    \]
    where $D_0$ is the total bond dimension of the all-zero blocks.
\end{theorem}

\begin{proof}\leanok
    \uses{thm:irred_decomp}
    Apply the irreducible block decomposition (Theorem~\ref{thm:irred_decomp})
    and partition the blocks by whether each has a nonzero Kraus operator; $D_0$
    accumulates the bond dimensions of the all-zero blocks. For every all-zero
    irreducible block $C_j$ and every $N \ge 1$, $V^{(N)}(C_j)_\sigma = 0$, so
    \[
        V^{(N)}(A)_\sigma
        = V^{(N)}\!(\textstyle\bigoplus_{k=1}^r B_k)_\sigma.
    \]
    At $N = 0$ the MPV coefficient is the bond dimension, giving
    $D = D_0 + \sum_{k=1}^r D_k$.
\end{proof}

\section{After-blocking decompositions and zero-tail identities}

\begin{theorem}[Blockwise TP-gauge normalization of irreducible blocks]%
    \label{thm:blockwise_tp_gauge}
    \lean{MPSTensor.exists_tp_gauge_blockwise}
    \leanok
    \uses{def:irreducible_tensor}
    Suppose $A$ generates the same MPV family as $\bigoplus_{k=1}^r B_k$, where
    each $B_k$ is irreducible with at least one nonzero Kraus operator. Then there
    are nonzero weights $(\mu_k)_{k=1}^r$ and blocks $(C_k)_{k=1}^r$ such that, for
    every $N$ and every $\sigma$,
    \[
        V^{(N)}(A)_\sigma
        =
        V^{(N)}\!(\textstyle\bigoplus_{k=1}^r \mu_k C_k)_\sigma
        =
        \sum_{k=1}^r \mu_k^{\,N} V^{(N)}(C_k)_\sigma,
    \]
    each $C_k$ irreducible and left-canonical,
    $\sum_i (C_k^i)^\dagger C_k^i = \Id$, with positive bond dimension. This is
    the blockwise form of the trace-preserving gauge; the arbitrary-tensor
    statement is Theorem~\ref{thm:tp_gauge_arbitrary}.
\end{theorem}

\begin{proof}\leanok
    The Perron--Frobenius TP-gauge construction is applied to each irreducible
    block. Gauge equivalence preserves irreducibility and the MPV family, and the
    positive scaling factors are absorbed into the weights $\mu_k$.
\end{proof}

\begin{theorem}[Primitive nonzero blocks are normal]%
    \label{thm:normal_live_block_primitive}
    \lean{MPSTensor.isNormal_live_block_of_primitive}
    \leanok
    \uses{def:normal, def:irreducible_tensor}
    Let $A$ be irreducible of positive bond dimension. If
    $\sum_i (A^i)^\dagger A^i = \Id$ and $\sigma_\partial(\E_A)=\{1\}$, then $A$ is
    normal. This is the nonzero-block instance of
    Theorem~\ref{thm:normal_tp_primitive_irred}.
\end{theorem}

\begin{proof}\leanok
    \uses{thm:normal_tp_primitive_irred}
    Apply Theorem~\ref{thm:normal_tp_primitive_irred}.
\end{proof}

\begin{theorem}[Blocking preserves irreducibility for primitive irreducible blocks]%
    \label{thm:blocked_irreducible_tp_primitive}
    \lean{MPSTensor.isIrreducibleTensor_blockTensor_of_tp_primitive_irr}
    \leanok
    \uses{def:blocked_tensor, def:irreducible_tensor}
    Let $A$ be left-canonical, irreducible, primitive, of positive bond
    dimension. For every $P>0$, the blocked tensor $A^{[P]}$ is irreducible:
    there is no orthogonal projection $Q$ with $Q \ne 0$, $Q \ne \Id$ such that,
    for every $P$-blocked index $\alpha$,
    \[
        (\Id-Q)(A^{[P]})^\alpha Q=0.
    \]
\end{theorem}

\begin{proof}\leanok
    A positive-definite fixed point of $\E_A$ remains fixed by the blocked transfer
    map; primitivity gives uniqueness of PSD fixed points for the blocked map,
    hence its transfer map is irreducible, which is the irreducibility criterion
    for $A^{[P]}$.
\end{proof}

\begin{theorem}[Extra blocking of primitive irreducible sector blocks]%
    \label{thm:tp_primitive_irreducible_extra_blocking}
    \lean{MPSTensor.tp_primitive_irreducible_extra_blocking}
    \leanok
    \uses{def:blocked_tensor, thm:blocked_irreducible_tp_primitive}
    Let $A$ be trace-preserving, primitive, and tensor-irreducible. For every
    $k>0$,
    \[
        \sum_\alpha ((A^{[k]})^\alpha)^\dagger
        (A^{[k]})^\alpha = \Id,
        \qquad
        \sigma_\partial(\E_{A^{[k]}})=\{1\},
    \]
    and $A^{[k]}$ is tensor-irreducible. This $k$ is a later
    common-refinement or injectivity length, distinct from the period-removal
    length that produced the sector block.
\end{theorem}

\begin{proof}\leanok
    \uses{thm:blocked_irreducible_tp_primitive}
    Trace preservation follows from blocking a trace-preserving tensor,
    primitivity from a positive power of a primitive transfer map, and
    irreducibility from Theorem~\ref{thm:blocked_irreducible_tp_primitive}.
\end{proof}

\begin{theorem}[Sorted TP primitive irreducible family forms a normal canonical form]%
    \label{thm:ncf_sorted_tp_primitive_irred}
    \lean{MPSTensor.isNormalCanonicalForm_of_tp_primitive_irr_sorted}
    \leanok
    \uses{def:is_normal_canonical_form}
    Let $(\mu_k,A_k)_{k=1}^r$ have $|\mu_1| \ge \cdots \ge |\mu_r|$ and
    $\mu_k \ne 0$ for every $k$. If every block is left-canonical, primitive,
    irreducible, and of positive bond dimension, then $(\mu_k,A_k)_{k=1}^r$
    satisfies the six conditions of Definition~\ref{def:is_normal_canonical_form}.
\end{theorem}

\begin{proof}\leanok
    \uses{def:is_normal_canonical_form}
    The hypotheses are exactly the constructor for normal canonical form.
\end{proof}

\begin{theorem}[Normal canonical form from an already primitive block decomposition]%
    \label{thm:ncf_from_primitive_input}
    \lean{MPSTensor.exists_normalCanonicalForm_of_primitive_input}
    \leanok
    \uses{def:is_normal_canonical_form}
    Suppose $A$ admits a block decomposition into irreducible left-canonical
    primitive blocks with pairwise distinct weight norms, nonzero weights, and
    positive bond dimensions. Then there exist $p>0$, weights $(\nu_k)$, and
    blocks $(B_k)$ satisfying Definition~\ref{def:is_normal_canonical_form} such
    that, for every blocked $N$ and $\sigma$,
    \[
        V^{(N)}\!(A^{[p]})_\sigma
        =
        V^{(N)}\!(\textstyle\bigoplus_k \nu_k B_k)_\sigma.
    \]
    This is the after-blocking form of
    Theorem~\ref{thm:exists_normal_canonical_form}; the distinct-modulus
    hypothesis is the subcase discussed in the remark after
    Definition~\ref{def:is_normal_canonical_form}.
\end{theorem}

\begin{proof}\leanok
    \uses{thm:exists_normal_canonical_form}
    Apply Theorem~\ref{thm:exists_normal_canonical_form} to the given primitive
    block decomposition.
\end{proof}

\begin{theorem}[Common blocking period for two primitive normal tensors]%
    \label{thm:common_period_two_tensors}
    \lean{MPSTensor.bilateral_commonPeriod_blocking_tp_primitive_normal}
    \leanok
    \uses{def:blocked_tensor, def:normal}
    Let $A$ and $B$ be left-canonical normal tensors of positive bond dimensions,
    with $A^{[p_A]}$, $B^{[p_B]}$ having primitive transfer maps for some
    $p_A,p_B>0$. Then there is a common $p>0$ such that both blocked tensors are
    left-canonical,
    \[
        \sum_\alpha ((A^{[p]})^\alpha)^\dagger
        (A^{[p]})^\alpha = \Id,
        \qquad
        \sum_\beta ((B^{[p]})^\beta)^\dagger
        (B^{[p]})^\beta = \Id,
    \]
    both have primitive transfer maps,
    $\sigma_\partial(\E_{A^{[p]}})=\{1\}$ and
    $\sigma_\partial(\E_{B^{[p]}})=\{1\}$, and both are normal.
\end{theorem}

\begin{proof}\leanok
    \uses{thm:primitive_blocking_divisible, thm:isNormal_blockTensor}
    Take $p = \operatorname{lcm}(p_A,p_B)$. Primitivity follows from
    divisibility, and left-canonical normalization and normality are preserved
    under blocking.
\end{proof}

\begin{remark}[Two-tensor primitive block decomposition]%
    \label{rem:ft_after_blocking_structural}
    Applying Theorem~\ref{thm:tp_primitive_blockdecomp} to each of two same-MPV
    tensors $A$ and $B$ produces, after a positive blocking length, a
    left-canonical primitive nonzero block decomposition for each side, with the
    separate zero-block count. This is a preliminary step toward
    Theorem~\ref{thm:ft_after_blocking_common_length_common_sector_theorem}, where
    the two sides are matched.
\end{remark}

\begin{theorem}[Cyclic sectors after the zero-block split]%
    \label{thm:ft_after_blocking_per_block_cyclic_live_zero_tail}
    \lean{MPSTensor.afterBlocking_perBlockCyclicData_of_sameMPV₂}
    \leanok
    \uses{thm:tp_gauge_arbitrary, def:primitive_irreducible_cyclic_sectors,
        thm:has_primitive_irreducible_cyclic_sectors_tp_irr}
    If $A$ and $B$ generate the same MPV family, then after the zero-block split
    and the trace-preserving gauge there are weighted nonzero-block tensors
    $\bigoplus_j \mu^A_j A_j$ and $\bigoplus_k \mu^B_k B_k$ with, for every
    $N\ge 1$ and $\sigma$,
    \[
        V^{(N)}(A)_\sigma
        =
        V^{(N)}\!(\textstyle\bigoplus_j \mu^A_j A_j)_\sigma,
        \qquad
        V^{(N)}(B)_\sigma
        =
        V^{(N)}\!(\textstyle\bigoplus_k \mu^B_k B_k)_\sigma,
    \]
    and the two nonzero parts agree:
    \[
        V^{(N)}\!(\textstyle\bigoplus_j \mu^A_j A_j)_\sigma
        =
        V^{(N)}\!(\textstyle\bigoplus_k \mu^B_k B_k)_\sigma.
    \]
    Each nonzero block has primitive irreducible cyclic sectors: for each $j$,
    there is $m_j\ge 1$ with a unit-weight sector decomposition
    $A_j^{[m_j]} \sim \bigoplus_{u=1}^{m_j} C^A_{j,u}$, and similarly for $B_k$.
\end{theorem}

\begin{proof}\leanok
    \uses{thm:tp_gauge_arbitrary,
        thm:has_primitive_irreducible_cyclic_sectors_tp_irr}
    Apply Theorem~\ref{thm:tp_gauge_arbitrary} separately to $A$ and $B$. Its
    positive-length conclusion gives, for $N\ge 1$,
    \[
        V^{(N)}(A)_\sigma
        =
        V^{(N)}\!(\textstyle\bigoplus_j \mu^A_j A_j)_\sigma,
        \qquad
        V^{(N)}(B)_\sigma
        =
        V^{(N)}\!(\textstyle\bigoplus_k \mu^B_k B_k)_\sigma.
    \]
    Chaining these through $V(A)=V(B)$ identifies the two nonzero parts. Finally
    apply Theorem~\ref{thm:has_primitive_irreducible_cyclic_sectors_tp_irr} to
    every trace-preserving irreducible nonzero-weight block.
\end{proof}

\begin{theorem}[Common blocking length for cyclic sectors]%
    \label{thm:ft_after_blocking_common_length_common_sector_theorem}
    \lean{MPSTensor.afterBlocking_commonLengthCommonSectorData_of_sameMPV₂}
    \leanok
    \uses{thm:ft_after_blocking_per_block_cyclic_live_zero_tail,
        thm:exists_common_blocked_cyclic_sector_family_common_multiple,
        thm:same_mpv2_pos_block_tensor_to_tensor_from_blocks,
        thm:same_mpv2_pos_to_tensor_from_blocks_block_power,
        thm:common_blocked_cyclic_sector_reindexed_samempv,
        thm:common_blocked_cyclic_sector_derived_properties}
    If two tensors generate the same MPV family, then the cyclic-sector
    decompositions on the two sides can be expressed at one common positive
    blocking length $p$. For every $N\ge 1$ and every blocked $\sigma$, each
    blocked tensor equals its powered nonzero part,
    \[
        V^{(N)}\!(A^{[p]})_\sigma
        =
        V^{(N)}\!(\textstyle\bigoplus_j (\mu^A_j)^p (A_j)^{[p]})_\sigma,
        \qquad
        V^{(N)}\!(B^{[p]})_\sigma
        =
        V^{(N)}\!(\textstyle\bigoplus_k (\mu^B_k)^p (B_k)^{[p]})_\sigma,
    \]
    and the two powered nonzero parts agree,
    \[
        V^{(N)}\!\bigl(\textstyle\bigoplus_j
        (\mu^A_j)^p (A_j)^{[p]}\bigr)_\sigma
        =
        V^{(N)}\!\bigl(\textstyle\bigoplus_k
        (\mu^B_k)^p (B_k)^{[p]}\bigr)_\sigma.
    \]
\end{theorem}

\begin{proof}\leanok
    \uses{thm:ft_after_blocking_per_block_cyclic_live_zero_tail,
        thm:exists_common_blocked_cyclic_sector_family_common_multiple,
        thm:same_mpv2_pos_block_tensor_to_tensor_from_blocks,
        thm:same_mpv2_pos_to_tensor_from_blocks_block_power,
        thm:common_blocked_cyclic_sector_reindexed_samempv,
        thm:common_blocked_cyclic_sector_derived_properties}
    Start from
    Theorem~\ref{thm:ft_after_blocking_per_block_cyclic_live_zero_tail}. With
    $p_A,p_B$ the least common multiples of the period-removal lengths on the two
    sides, use $p=\operatorname{lcm}(p_A,p_B)$. The prescribed-common-multiple
    lemma constructs both one-sided sector families at this length, and the
    positive-length blocking lemmas carry the one-sided tensor/nonzero-part
    equalities over to the common alphabet; for every $N\ge 1$,
    \[
        V^{(N)}\!(A^{[p]})_\sigma
        =
        V^{(N)}\!\bigl(\textstyle\bigoplus_j
        (\mu^A_j)^p (A_j)^{[p]}\bigr)_\sigma.
    \]
    The same applies on the $B$ side, and the two-sided blocking lemma gives
    equality of the two powered nonzero parts.
\end{proof}

\begin{theorem}[Primitive block decompositions at a common length]%
    \label{thm:after_blocking_common_primitive_irreducible_blocks_reindexed}
    \lean{MPSTensor.afterBlocking_commonPrimitiveIrreducibleBlocks_of_reindexedNonzeroParts}
    \leanok
    \uses{thm:ft_after_blocking_common_length_common_sector_theorem,
        thm:common_blocked_cyclic_sector_reindexed_nonzero_part,
        def:same_mpv2_pos}
    Assume the nonzero part can be identified with the common blocked alphabet.
    Then two tensors with the same MPV family admit one $p>0$ and nonzero
    weighted block families such that, for every $N\ge 1$,
    \[
        V^{(N)}\!(A^{[p]})_\sigma
        =
        V^{(N)}\!(\textstyle\bigoplus_x \mu^A_x A_x)_\sigma,
        \qquad
        V^{(N)}\!(B^{[p]})_\sigma
        =
        V^{(N)}\!(\textstyle\bigoplus_y \mu^B_y B_y)_\sigma,
    \]
    and the two nonzero parts agree at every positive length,
    \[
        V^{(N)}\!(\textstyle\bigoplus_x \mu^A_x A_x)_\sigma
        =
        V^{(N)}\!(\textstyle\bigoplus_y \mu^B_y B_y)_\sigma.
    \]
    The all-zero leftover is restored by the separate zero-tail cancellation
    theorem. All blocks are trace-preserving, primitive, tensor-irreducible, and
    of positive bond dimension, and all weights are nonzero.
\end{theorem}

\begin{proof}\leanok
    \uses{thm:ft_after_blocking_common_length_common_sector_theorem,
        thm:common_blocked_cyclic_sector_reindexed_nonzero_part}
    Apply
    Theorem~\ref{thm:ft_after_blocking_common_length_common_sector_theorem} to
    choose the common blocking length and the two common cyclic-sector families.
    The assumed word identification converts the blocked nonzero part into the
    weighted common-sector family on each side, and since the zero-tail term
    vanishes at positive length, this gives the displayed identities. The
    structural properties and nonzero weights come from the common-sector
    families.
\end{proof}

\section{Primitive-block decomposition of an arbitrary tensor}%
\label{sec:ch11b_primitive_block_reduction}

This section gives the primitive sector theorem used in the period-removal step
and the resulting primitive-block decomposition of an arbitrary tensor after
blocking.

\begin{theorem}[Unconditional primitive blocked sectors]%
    \label{thm:blocked_sector_unconditional}
    \lean{MPSTensor.primitive_and_irreducible_sectorBlocks_of_cyclic_decomp_after_blocking}
    \leanok
    \uses{thm:cyclic_sector_decomp_after_blocking,
        thm:sector_irred_unconditional,
        thm:primitive_adjoint_equiv}
    Let $A$ be a left-canonical irreducible tensor whose adjoint transfer map has
    peripheral spectrum $\{\gamma^j : 0 \le j < m\}$ for a primitive $m$th root of
    unity $\gamma$, with blocked cyclic-sector decomposition $(C_u)_{u=1}^m$,
    $(P_u)_{u=1}^m$, and
    $\varphi_u : \MN{\dim C_u} \xrightarrow{\sim} P_u \MN{D} P_u$. Then every
    sector tensor satisfies $\sigma_\partial(\E_{C_u})=\{1\}$ and is
    tensor-irreducible.
\end{theorem}

\begin{proof}\leanok
    \uses{thm:sector_irred_unconditional,
        thm:primitive_adjoint_equiv}
    Theorem~\ref{thm:sector_irred_unconditional} gives irreducibility of
    $(\E_A^\dagger)^m$ on each corner $P_u$. The compression equivalences
    $\varphi_u$ intertwine the adjoint transfer map of $C_u$ with the
    corresponding corner restriction of $(\E_A^\dagger)^m$, and the cyclic
    peripheral-spectrum argument shows each corner restriction is primitive.
    Primitivity and irreducibility are preserved across $\varphi_u$; then use
    Theorem~\ref{thm:primitive_adjoint_equiv} to pass from the adjoint blocked
    map to the primal transfer map of $C_u$.
\end{proof}

\begin{theorem}[Arbitrary tensor to primitive block decomposition]%
    \label{thm:tp_primitive_blockdecomp}
    \lean{MPSTensor.exists_tp_primitive_blockDecomp_after_blocking}
    \leanok
    \uses{def:blocked_tensor, def:same_mpv2_pos}
    For any MPS tensor $A$, there exist $p\ge 1$, a zero-tail bond dimension
    $D_0$, nonzero weights $(\mu_k)_{k=1}^r$, and blocks $(B_k)_{k=1}^r$ such
    that, for every blocked $\sigma$ of positive length $N$,
    \[
        V^{(N)}\!(A^{[p]})_\sigma
        = V^{(N)}\!(\textstyle\bigoplus_{k=1}^{r} \mu_k B_k)_\sigma.
    \]
    The zero-block contribution is recorded separately by the length-zero
    identity $D_0 + \sum_{k=1}^{r} D_k = D$. Each block is left-canonical and
    primitive,
    \[
        \sum_i (B_k^i)^\dagger B_k^i = \Id,
        \qquad
        \sigma_\partial(\E_{B_k})=\{1\},
    \]
    and each bond dimension is positive.
\end{theorem}

\begin{proof}\leanok
    \uses{thm:tp_gauge_arbitrary, thm:common_blocking_period,
        lem:block_weights_ne_zero,
        thm:same_mpv2_pos_block_tensor_to_tensor_from_blocks}
    Theorem~\ref{thm:tp_gauge_arbitrary} supplies the zero-block dimension, the
    left-canonical nontrivial blocks, the positive-length equality, and the
    identity $D = D_0 + \sum_k D_k$.
    Theorem~\ref{thm:common_blocking_period} provides a common blocking period
    $p$ making all transfer maps primitive,
    Lemma~\ref{lem:block_weights_ne_zero} keeps the blocked weights nonzero, and
    Lemma~\ref{thm:same_mpv2_pos_block_tensor_to_tensor_from_blocks} carries
    the positive-length equality through blocking.
\end{proof}

\begin{theorem}[Two tensors after one-sided primitive reductions]%
    \label{thm:after_blocking_tp_primitive_blockdecompositions_same_mpv2}
    \lean{MPSTensor.afterBlocking_tpPrimitiveBlockDecompositions_of_sameMPV₂}
    \leanok
    \uses{thm:tp_primitive_blockdecomp, def:same_mpv2, def:same_mpv2_pos}
    If $A$ and $B$ have the same MPV family, then applying
    Theorem~\ref{thm:tp_primitive_blockdecomp} separately to $A$ and $B$ gives
    positive blocking lengths $p_A,p_B$ and primitive weighted block sums
    $A_{\mathrm{nz}}$, $B_{\mathrm{nz}}$ such that $A^{[p_A]}$ and $B^{[p_A]}$
    have the same MPV family, $A^{[p_B]}$ and $B^{[p_B]}$ have the same MPV
    family, $A^{[p_A]}$ agrees with $A_{\mathrm{nz}}$ at every positive length,
    and $B^{[p_B]}$ agrees with $B_{\mathrm{nz}}$ at every positive length. All
    displayed nonzero blocks are trace-preserving and primitive, of positive bond
    dimension, with nonzero weights.
\end{theorem}

\begin{proof}\leanok
    \uses{thm:tp_primitive_blockdecomp, def:same_mpv2, def:same_mpv2_pos}
    Apply Theorem~\ref{thm:tp_primitive_blockdecomp} to each tensor. Equality of
    MPV families is preserved by blocking, and the all-zero summand has zero MPV
    coefficients at positive length, so each blocked tensor agrees there with its
    nonzero weighted block sum.
\end{proof}
